% Canonical Paper II source.
The intrinsic complex structure of a regular AES makes holomorphic pullback a
natural source of new regular examples.  This section records the construction in
the form needed by Paper III.  Every statement below is confined to the locus on
which the relevant holomorphic map is a local biholomorphism.  No assertion is made
about extending through a critical or branch set, or through a zero or pole of the
cylindrical target map.

\subsection{A planar harmonic base AES}

Fix \(\phi\in\R\), write \(w=\xi+i\eta\), and define
\begin{equation}\label{eq:planar-base-assignment}
 A_\phi(w):=\operatorname{Im}(e^{-i\phi}w),
 \qquad
 Q_\phi(w)^2:=\mu^2+\lambda^2A_\phi(w)^2,
\end{equation}
where \(\mu\ne0\) and \(\lambda\in\R\).  On \(\C\), set
\begin{equation}\label{eq:planar-base-metric}
 h_\phi:=\frac{|dw|^2}{Q_\phi(w)^2}.
\end{equation}
The phase merely rotates the Euclidean coordinate in which the assignment is the
second coordinate, but retaining it makes the zero direction explicit.

\begin{proposition}[Planar harmonic AES and curvature law]
\label{prop:planar-harmonic-aes}
The tuple \((\C,h_\phi,A_\phi;\mu,\lambda)\) is a regular AES and
\begin{align}
 |\nabla A_\phi|_{h_\phi}^2
 &=\mu^2+\lambda^2A_\phi^2,
 \label{eq:planar-base-eikonal}\\
 \Delta_{h_\phi}A_\phi&=0,
 \label{eq:planar-base-harmonicity}\\
 K_{h_\phi}
 &=\lambda^2
   \frac{\mu^2-\lambda^2A_\phi^2}
        {\mu^2+\lambda^2A_\phi^2}.
 \label{eq:planar-base-curvature}
\end{align}
In particular, the AES equation by itself neither determines the curvature nor
forces the assignment to satisfy the pointwise equation
\(\Delta_g a=2\lambda^2a\) of the basic hyperbolic model.
\end{proposition}

\begin{proof}
After the Euclidean rotation \(u+i v=e^{-i\phi}w\), one has
\(A_\phi=v\) and
\[
 h_\phi=e^{2\rho(v)}(du^2+dv^2),
 \qquad
 \rho(v)=-\frac12\log(\mu^2+\lambda^2v^2).
\]
Since the inverse metric is \(e^{-2\rho}\) times the Euclidean inverse metric,
\[
 |\nabla v|_{h_\phi}^2=e^{-2\rho}
 =\mu^2+\lambda^2v^2.
\]
In real dimension two the scalar Laplacian obeys
\(\Delta_{e^{2\rho}g_0}f=e^{-2\rho}\Delta_{g_0}f\); hence
\(\Delta_{h_\phi}v=0\).  Finally,
\[
 \rho''(v)
 =-\lambda^2\frac{\mu^2-\lambda^2v^2}
                         {(\mu^2+\lambda^2v^2)^2},
\]
and the conformal curvature formula
\(K=-e^{-2\rho}\Delta_{g_0}\rho\) gives
\eqref{eq:planar-base-curvature}.
\end{proof}

The zero locus is the Euclidean line
\begin{equation}\label{eq:planar-base-zero-line}
 Z(A_\phi)=e^{i\phi}\R.
\end{equation}
It is regular because \(|\nabla A_\phi|_{h_\phi}=|\mu|\) there.  The next theorem
shows that its inverse images give all zero curves produced by this base model on
the regular locus.

\subsection{Locally biholomorphic pullback}

\begin{theorem}[Holomorphic pullback of the planar base]
\label{thm:planar-holomorphic-pullback}
Let \(\Sigma\) be a Riemann surface, \(U\subset\Sigma\) open, and
\(F:U\to\C\) holomorphic with \(dF_p\ne0\) for every \(p\in U\).  Define
\begin{equation}\label{eq:planar-pullback-data}
 a_F:=A_\phi\circ F,
 \qquad
 g_F:=F^*h_\phi.
\end{equation}
Then \((U,g_F,a_F;\mu,\lambda)\) is a regular AES.  Moreover,
\begin{align}
 \Delta_{g_F}a_F&=0,
 \label{eq:planar-pullback-harmonicity}\\
 K_{g_F}
 &=\lambda^2
   \frac{\mu^2-\lambda^2a_F^2}
        {\mu^2+\lambda^2a_F^2}.
 \label{eq:planar-pullback-curvature}
\end{align}
If \(z\) is a local complex coordinate on \(U\), then
\begin{equation}\label{eq:planar-pullback-local-form}
 g_F=
 \frac{|F'(z)|^2}{\mu^2+\lambda^2a_F(z)^2}|dz|^2,
 \qquad
 Z(a_F)=F^{-1}(e^{i\phi}\R).
\end{equation}
Finally, every holomorphic field \(\Psi\) on an open neighborhood of \(F(U)\)
pulls back to the arithmetic holomorphic field \(\Psi\circ F\) on \(U\).
\end{theorem}

\begin{proof}
The nonvanishing of \(dF\) makes \(F:(U,g_F)\to(\C,h_\phi)\) a local oriented
isometry.  The norm of the gradient, the scalar Laplacian, and Gaussian curvature
are preserved under local isometries, so
Proposition~\ref{prop:planar-harmonic-aes} gives the AES equation and
\eqref{eq:planar-pullback-harmonicity}--\eqref{eq:planar-pullback-curvature}.
The coordinate expression follows directly from
\(F^*|dw|^2=|F'(z)|^2|dz|^2\), and the zero-set identity follows from
\eqref{eq:planar-base-zero-line}.  Since \(F\) is holomorphic for the given
Riemann-surface structures, so is \(\Psi\circ F\); by
Definition~\ref{def:surface-arithmetic-holomorphic}, this is arithmetic
holomorphicity for \(g_F\).
\end{proof}

\begin{warning}[The critical set is not part of the theorem]
\label{warn:planar-pullback-critical-set}
If \(F'(p)=0\), formula~\eqref{eq:planar-pullback-local-form} defines a degenerate
quadratic form at \(p\), not a Riemannian metric.  If, in addition,
\(A_\phi(F(p))=0\), the inverse image of the target zero line may have several
local branches.  Such a point cannot be regular: the AES equation would require
\(|\nabla a_F|=|\mu|\) on its zero set.  Branched pullbacks and their zero-graph
normal forms belong to the singular theory of Paper III.
\end{warning}

\subsection{The complete cylindrical AES}

The planar construction is adapted to inverse images of a line.  For inverse images
of a circle, the appropriate regular base is the logarithmic cylinder.  Let
\(W\) denote the standard coordinate on \(\C^*=\C\setminus\{0\}\), and put
\begin{equation}\label{eq:cylinder-s-coordinate}
 s(W):=\log|W|,
 \qquad
 g_{\mathrm{cyl}}:=\left|\frac{dW}{W}\right|^2.
\end{equation}
Define
\begin{equation}\label{eq:cylinder-assignment}
 a_{\mathrm{cyl}}(W):=
 \begin{cases}
  \dfrac{\mu}{\lambda}\sinh(\lambda s(W)),&\lambda\ne0,\\[6pt]
  \mu s(W),&\lambda=0.
 \end{cases}
\end{equation}

\begin{proposition}[Cylindrical AES]
\label{prop:cylindrical-aes}
The tuple
\((\C^*,g_{\mathrm{cyl}},a_{\mathrm{cyl}};\mu,\lambda)\) is a complete regular
AES.  It is flat and satisfies
\begin{align}
 |\nabla a_{\mathrm{cyl}}|_{g_{\mathrm{cyl}}}^2
 &=\mu^2+\lambda^2a_{\mathrm{cyl}}^2,
 \label{eq:cylinder-eikonal}\\
 \Delta_{g_{\mathrm{cyl}}}a_{\mathrm{cyl}}
 &=\lambda^2a_{\mathrm{cyl}},
 \label{eq:cylinder-laplacian}\\
 Z(a_{\mathrm{cyl}})&=\{|W|=1\}.
 \label{eq:cylinder-zero-circle}
\end{align}
For \(\lambda=0\), equation~\eqref{eq:cylinder-laplacian} has the stated limiting
meaning \(\Delta a_{\mathrm{cyl}}=0\).
\end{proposition}

\begin{proof}
Locally choose a logarithm \(\zeta=\log W=s+i\theta\).  Then
\[
 g_{\mathrm{cyl}}=|d\zeta|^2=ds^2+d\theta^2.
\]
The universal covering map identifies this metric with the quotient of the
Euclidean plane by \(\theta\mapsto\theta+2\pi\), so it is complete and flat.  For
\(\lambda\ne0\),
\[
 \frac{da_{\mathrm{cyl}}}{ds}=\mu\cosh(\lambda s),
 \qquad
 \frac{d^2a_{\mathrm{cyl}}}{ds^2}
 =\lambda^2a_{\mathrm{cyl}}.
\]
The identity \(\cosh^2t=1+\sinh^2t\) proves
\eqref{eq:cylinder-eikonal}, while the flat-cylinder Laplacian proves
\eqref{eq:cylinder-laplacian}.  The formulas for \(\lambda=0\) follow directly.
Since \(\sinh t=0\) only at \(t=0\), the zero set is \(s=0\), equivalently
\(|W|=1\).
\end{proof}

As in Theorem~\ref{thm:planar-holomorphic-pullback}, any holomorphic local
biholomorphism \(W:U\to\C^*\) gives a regular pullback AES
\begin{equation}\label{eq:cylinder-pullback-data}
 g_W=\left|\frac{dW}{W}\right|^2,
 \qquad
 a_W=
 \begin{cases}
  \dfrac{\mu}{\lambda}\sinh(\lambda\log|W|),&\lambda\ne0,\\[6pt]
  \mu\log|W|,&\lambda=0,
 \end{cases}
\end{equation}
with regular zero locus
\begin{equation}\label{eq:cylinder-pullback-zero-set}
 Z(a_W)=W^{-1}(S^1).
\end{equation}
On this pullback one has
\(K_{g_W}=0\) and \(\Delta_{g_W}a_W=\lambda^2a_W\).

\subsection{Automorphic descent and the Paper III interface}

The preceding formula is useful when \(W\) is constructed from an automorphic or
meromorphic function.  The following elementary descent statement isolates what is
already regular analysis from what remains singular geometry.

\begin{proposition}[Unitary automorphy and descent]
\label{prop:cylinder-automorphic-descent}
Let a group \(\Gamma\) act biholomorphically on a Riemann surface \(U\), and let
\(W:U\to\C^*\) be a holomorphic local biholomorphism.  Suppose there is a character
\(\chi:\Gamma\to S^1\) such that
\begin{equation}\label{eq:cylinder-unitary-automorphy}
 W(\gamma p)=\chi(\gamma)W(p).
\end{equation}
Then \(g_W\), \(a_W\), and \(Z(a_W)\) are \(\Gamma\)-invariant and descend to every
smooth quotient on which the action is free and properly discontinuous.

More generally, suppose the target transformations define an action and have the
form
\begin{equation}\label{eq:cylinder-dihedral-automorphy}
 W(\gamma p)=\chi_\gamma W(p)^{\varepsilon(\gamma)},
 \qquad
 |\chi_\gamma|=1,
 \qquad
 \varepsilon(\gamma)\in\{1,-1\}.
\end{equation}
Then \(\varepsilon:\Gamma\to\{1,-1\}\) is a homomorphism, \(g_W\) and
\(Z(a_W)\) are invariant, and \(a_W\) is multiplied by \(\varepsilon(\gamma)\).
It therefore descends as a scalar only on the sign-preserving subgroup; on the full
quotient it is naturally a section of the corresponding real sign line bundle.
\end{proposition}

\begin{proof}
Multiplication by a unit complex number preserves \(\log|W|\) and \(dW/W\).
Inversion sends these quantities to \(-\log|W|\) and \(-dW/W\), respectively.
Thus the metric and the unit-circle preimage are unchanged, whereas the odd
function in~\eqref{eq:cylinder-assignment} changes sign.  Standard quotient descent
then proves the statement.
\end{proof}

\begin{remark}[Converting an automorphic interval into the zero circle]
\label{rem:automorphic-interval-cayley}
Let \(\beta\) be a meromorphic automorphic function.  On an open sheet avoiding its
branch values, choose \(R\) and define \(W\) by
\begin{equation}\label{eq:automorphic-cayley-interface}
 R^2=\frac{\beta}{1-\beta},
 \qquad
 W=\frac{R-i}{R+i}.
\end{equation}
Where these expressions are finite and locally biholomorphic,
\begin{equation}\label{eq:automorphic-interval-zero-equivalence}
 |W|=1
 \quad\Longleftrightarrow\quad
 R\in\R\cup\{\infty\}
 \quad\Longleftrightarrow\quad
 \beta\in[0,1].
\end{equation}
Changing the square-root sheet sends \(R\) to \(-R\) and \(W\) to \(W^{-1}\).
Thus the cylindrical metric and zero locus are independent of that sheet change,
whereas the assignment changes sign.  This observation supplies the exact regular
analytic interface for a normalized Hauptmodul or similar automorphic coordinate.
It does not assert that \(\beta^{-1}([0,1])\) has any particular global graph type;
that conclusion requires separate ramification and tessellation data.
\end{remark}

For a meromorphic or algebraic automorphic construction, define its Paper II
regular locus to be the union of open sheets on which
\begin{equation}\label{eq:automorphic-regular-locus}
 0<|W|<\infty,
 \qquad
 dW\ne0,
\end{equation}
and on which a required algebraic branch has been chosen.  Equations
\eqref{eq:cylinder-pullback-data}--\eqref{eq:cylinder-pullback-zero-set} apply on
that locus without change.  Zeros and poles of \(W\) are excluded because
\(\log|W|\) is not a finite smooth assignment there.  Critical points are excluded
because the pulled-back metric degenerates, and algebraic branch points require a
separate local uniformization.  These excluded sets are among the natural loci at
which inverse images of \(S^1\) may meet or branch, and at which they may
accumulate.
Boundary escape and nonproperness can also cause accumulation without appearing
in this local list.  Paper II exports the regular formula and the descent criterion;
Paper III owns the completion, local branching, singular zero graph, and
parameterized tube topology.
